% =====================================================================
%  Chapter 2 — State of the Art
%  All 16 cite keys resolve against biblio_thesis.bib (verified).
%  Style notes: key terms in \textbf so the chapter can be skimmed;
%  enumerations use itemize rather than running prose.
% =====================================================================

This chapter reviews the literature at the intersection of contactless radar
sensing, micro-Doppler signal processing and machine-learning gait analysis, and
positions the present work within it. The survey is organised from the physical
substrate upward:

\begin{itemize}
  \item The case for radar as a gait sensor (Section~\ref{sec:contactless});
  \item What a radar micro-Doppler signature physically encodes, and how reliably
        it can be measured (Section~\ref{sec:foundations});
  \item How deep learning has been applied to micro-Doppler signatures
        (Section~\ref{sec:dl});
  \item Strategies for learning from the small datasets typical of this field
        (Section~\ref{sec:small});
  \item The emerging use of radar for neurodegenerative disease
        (Section~\ref{sec:neuro});
  \item The validation pitfalls that recur throughout
        (Section~\ref{sec:pitfalls}).
\end{itemize}

\noindent Section~\ref{sec:keyworks} then condenses the most influential
works into Table~\ref{tab:keyworks}, and Section~\ref{sec:gap} distils the
resulting research gap that this thesis addresses.

\section{Contactless gait sensing}
\label{sec:contactless}

Quantitative gait analysis is clinically valuable for detecting and monitoring
movement disorders, but each of the dominant instrumented modalities carries a
cost. \textbf{Wearable inertial sensors} are accurate and inexpensive, yet they
must be worn, which constrains continuous use and biases natural gait.
\textbf{Optical motion capture} and video are accurate but expensive,
illumination-dependent and privacy-invasive.

\textbf{Radar} has emerged as a complementary alternative that is
\textbf{contactless}, insensitive to lighting and clothing, and
\textbf{privacy-preserving}, which makes it attractive for continuous and in-home
monitoring \cite{seifert2020,gurbuz2024,alanazi2022}. These properties motivate
the growing body of work that uses radar to characterise human gait, reviewed in
the remainder of this chapter.

\section{Micro-Doppler foundations and gait validation}
\label{sec:foundations}

When a radar illuminates a walking person, different body parts moving at
different radial velocities induce distinct Doppler shifts. The resulting
time--frequency representation, the \textbf{micro-Doppler signature}, encodes the
slow, steady motion of the \textbf{torso} at low Doppler frequencies and the
faster swings of the \textbf{limbs} at higher frequencies \cite{seifert2020}.

Establishing which parts of this signature are physically trustworthy has been
the subject of careful validation studies. Seifert et al.\ \cite{seifert2020}
benchmarked a 24\,GHz Doppler radar against marker-based motion capture and
instrumented-treadmill ground-reaction forces. They showed that radar recovers
\textbf{rhythm and pace} parameters (stride and step time, cadence, stride
length) reliably, but that joint-velocity estimates and gait-event timing depend
strongly on \textbf{geometry, viewing angle and walking speed}.

Complementing this, the survey of Gurbuz et al.\ \cite{gurbuz2024} catalogues
radar-based gait-parameter estimation for fall-risk assessment and notes a
crucial limitation: radar captures cadence and step time well, but measures
\textbf{gait variability} only poorly to moderately. This matters because
stride-to-stride fluctuation is a defining sign of parkinsonian gait.

\subsection{The source of the dataset used in this thesis}
\label{sec:dataset-origin}

The radar network that produced the dataset analysed in this work is that of
\textbf{López-Delgado et al.} \cite{lopezdelgado2026}, developed at the
Universidad Politécnica de Madrid in collaboration with the Hospital General
Universitario Gregorio Marañón de Madrid. That study sits squarely in the validation
tradition described above: it compares multi-node radar configurations against a
\textbf{Vicon} optical motion-capture reference \cite{merriaux2017} and
demonstrates accurate spatiotemporal
parameter extraction in both healthy and Parkinson's populations.

Two of its findings directly motivate design decisions taken later in this
thesis:

\begin{itemize}
  \item \textbf{Foot-oriented radar nodes best recover temporal gait parameters}
        in motorically impaired subjects. This justifies treating the foot and
        torso channels as distinct, complementary information sources rather than
        merging them into one.
  \item \textbf{The torso-velocity/heel-strike coupling breaks down in
        Parkinson's disease}, although it is reliable in healthy gait. This is a
        \emph{relational} property of the signature: a two-channel image
        model can represent it, whereas scalar summary features cannot.
\end{itemize}

Critically, however, that work stops at parameter \emph{measurement}: it contains
\textbf{no disease classifier}. This is the point of departure for the present
thesis, which takes the same acquisition system and asks whether the diagnosis
itself can be predicted directly from the micro-Doppler signatures.

\section[Deep learning for micro-Doppler gait classification]{Deep learning for micro-Doppler gait\\ classification}
\label{sec:dl}

A second line of work treats the micro-Doppler signature as an image or a
sequence, and applies machine learning to classify gait. Two architectural
families dominate:

\begin{itemize}
  \item \textbf{Convolutional networks (CNNs)} applied to the full
        time--frequency image.
  \item \textbf{Recurrent networks}, such as the \textbf{LSTM}, applied to
        one-dimensional velocity time-series, for example the upper, lower
        and power-weighted Doppler envelopes.
\end{itemize}

Hayashi et al.\ \cite{hayashi2021} exemplify the second family, classifying
young versus elderly gait from 24\,GHz micro-Doppler with a
\textbf{long short-term memory} network on extracted envelopes. They reach
\textbf{94.9\,\%} accuracy and report that the stance-phase (lower-envelope)
feature alone was most discriminative. A companion simulation study
\cite{hoshiga2021} corroborates that \textbf{leg-motion features} carry most of
the age-discriminative signal.

CNN-based approaches on micro-Doppler images are equally common across
human-activity and gait recognition, with reported accuracies frequently in
the 90--98\,\% range: Nguyen et al.\ \cite{nguyen2024} separate five gait
patterns at 96 to 99\,\% with a millimetre-wave sensor, Ha et al.\
\cite{ha2023} reach 97.9\,\% with a compact CNN on cepstral features from six
subjects, and Fard et al.\ \cite{fard2026} report 89.7\,\% under a
leave-one-person-out protocol on three subjects. Notably, Papanastasiou et
al.\ \cite{papanastasiou2020} and Ni and Huang \cite{ni2020} show that
micro-Doppler encodes \textbf{individual identity} strongly enough to recognise
people ($>$93.5\,\% and 96.7\,\% respectively). As Section~\ref{sec:pitfalls}
argues, that result has direct and cautionary implications for how disease
classifiers must be validated. Alternative front-ends have also been explored,
including audio-processing pipelines applied to the Doppler signature
\cite{ha2023} and low-cost millimetre-wave sensing for gait \cite{alanazi2022}.

\section{Learning from small radar datasets}
\label{sec:small}

Radar gait datasets are typically small, on the order of tens of subjects,
because data collection requires supervised sessions with specialised hardware.
Two strategies recur for coping with this scarcity:

\begin{itemize}
  \item \textbf{Transfer learning.} Seyfio\u{g}lu et al.\ \cite{seyfioglu2018}
        show that deep networks pre-trained on diversified micro-Doppler transfer
        effectively to new motion-classification tasks (97\,\% on a
        seven-class task against 86\,\% for the same network trained from
        scratch), and Park et al.\ \cite{park2016} demonstrated transfer
        learning on micro-Doppler as early as 2016, reaching 80.3\,\% on five
        subjects.
  \item \textbf{Synthetic data augmentation.} Erol et al.\ \cite{erol2020} use a
        kinematically sifted auxiliary-classifier \textbf{GAN} to synthesise
        realistic radar micro-Doppler signatures and expand small training sets,
        raising cross-environment accuracy from 84 to 93\,\%.
\end{itemize}

Both strategies are directly relevant to a \textbf{58-subject} study, and they
inform the transfer-learning and augmentation choices adopted in this thesis.

\section[Radar for neurodegenerative disease and Parkinson's]{Radar for neurodegenerative disease\\ and Parkinson's}
\label{sec:neuro}

Applying radar gait analysis specifically to neurodegenerative disease is recent
and still limited. The \textbf{STRIDE} system of Cai et al.\ \cite{cai2023} uses radar
intelligence for Alzheimer's-disease-and-related-dementia risk evaluation,
illustrating radar's promise for passive cognitive-decline screening.

For Parkinson's disease specifically, machine-learning gait classification has
been pursued mainly through \textbf{non-radar} modalities such as force sensors
and wearables: Meral and Ozbilgin \cite{neurosurgery2025} report 99.4\,\%
accuracy for Parkinson's versus control from force-plate recordings under
$k$-fold validation with no subject-independent test. The radar literature on Parkinson's
remains largely confined to parameter validation \cite{lopezdelgado2026} rather
than automated diagnosis. To the best of our knowledge, \textbf{no published work
performs subject-independent Parkinson's-versus-control classification directly
from radar micro-Doppler signatures}. That is the specific gap this thesis
targets. The accuracies reported across these works, gathered in
Table~\ref{tab:keyworks}, cluster between 90 and 99\,\%; the next section
explains why they must be read with caution.

\section{Limitations of the state of the art in the field}
\label{sec:pitfalls}

The high accuracies reported across the micro-Doppler classification literature
must be read against a set of recurring methodological hazards, because several
of them are known to inflate results. Four are relevant here.

\begin{itemize}
  \item \textbf{Subject leakage.} Because micro-Doppler encodes individual
        identity so strongly \cite{papanastasiou2020,ni2020}, any evaluation that
        allows recordings from the same subject to appear in both training and
        test sets permits the model to recognise \emph{people} rather than the
        target condition. \textbf{Subject-independent
        (leave-one-subject-out)} validation is therefore essential, yet much of
        the field relies on random hold-out splits
        \cite{hayashi2021,hoshiga2021}. Of the sixteen works reviewed in this
        chapter, \textbf{only two} report a genuinely subject-independent
        protocol \cite{fard2026,park2016}, and neither addresses Parkinson's
        disease. Put differently, a walk varies in two ways: within a person
        from one recording to the next (intra-patient variability) and
        between people (inter-patient variability). A model evaluated on
        people it has already seen only has to master the first; the systems
        in this literature are never shown to model the second, which is why
        they fail on new patients.

  \item \textbf{Shortcut learning on artifacts.} CNNs trained on micro-Doppler
        images often latch onto irrelevant signals rather than gait. Hayashi et al.\
        \cite{hayashi2021} explicitly demonstrate that an AlexNet reaching
        \textbf{97.8\,\%} did so by exploiting \textbf{background-noise
        differences}, which arose because their two groups were recorded at
        different locations. This failure mode is called \textbf{shortcut learning}: the network
        found an easier giveaway in the images, the noise pattern of each
        recording site, and used it instead of the gait. The same risk exists
        for any network trained on these images, which is why this thesis inspects, with
        attention maps, what its own networks actually look at.

  \item \textbf{The age confound.} Group-level confounds run through this whole literature, and
        \textbf{age} is especially dangerous for a Parkinson's study because
        it is so easy to miss. Because
        micro-Doppler alone separates young from elderly walkers at
        \textbf{94.9\,\%} \cite{hayashi2021}, and Parkinson's cohorts are
        typically older than controls, a classifier can appear to detect disease
        while in fact detecting age.

  \item \textbf{Measuring what matters least reliably.} Among the features
        most diagnostic of parkinsonian gait, increased
        \textbf{stride-to-stride variability} is exactly what radar measures
        least reliably \cite{gurbuz2024}, and joint-velocity estimates depend
        strongly on geometry \cite{seifert2020}. Pipelines
        that aggressively normalise amplitude, or that use very short temporal
        windows, therefore risk discarding the very pathology they seek.
\end{itemize}

\noindent Together these observations argue that \textbf{careful evaluation design} is
the central unsolved problem in this niche, rather than raw accuracy. Concretely,
that means subject-independent validation, explicit confound control, and
interpretability checks confirming the model attends to physiologically plausible
Doppler bands.

\section{Summary of the most influential works}
\label{sec:keyworks}

Table~\ref{tab:keyworks} condenses the studies that most directly shaped the
design of this thesis, pairing each key finding with the concrete methodological
consequence drawn from it.

\begin{table}[p]
\centering
\footnotesize
\caption[Key related works and their influence on this thesis]{Key related works,
their principal findings, and the specific methodological decision each one
motivates in this thesis.}
\label{tab:keyworks}
\begin{tabular}{@{}p{3.1cm}p{5.3cm}p{5.3cm}@{}}
\toprule
\textbf{Work} & \textbf{Key finding} & \textbf{Influence on this thesis} \\
\midrule

López-Delgado et al.\ \cite{lopezdelgado2026} \newline \emph{(UPM, source of our dataset)}
& \textbullet~Multi-node FMCW radar validated against Vicon. \newline \textbullet~Foot nodes best recover temporal parameters. \newline \textbullet~Contains \textbf{no classifier}.
& \textbullet~Provides the dataset and the geometry. \newline \textbullet~Motivates the \textbf{two-channel} input. \newline \textbullet~Defines the gap this work fills. \\
\addlinespace

Hayashi et al.\ \cite{hayashi2021}
& \textbullet~LSTM on Doppler envelopes: \textbf{94.9\,\%} young vs elderly. \newline \textbullet~Their AlexNet hit 97.8\,\% via a \textbf{site artifact}, not gait.
& \textbullet~\textbf{Age} treated as a dominant confound. \newline \textbullet~Attention checks made mandatory. \newline \textbullet~The strongest cautionary result. \\
\addlinespace

Gurbuz et al.\ \cite{gurbuz2024}
& \textbullet~Radar measures cadence and step time well. \newline \textbullet~Gait variability only poorly to moderately.
& \textbullet~Tempers expectations for Parkinson's. \newline \textbullet~Argues against aggressive amplitude normalisation. \\
\addlinespace

Seifert et al.\ \cite{seifert2020}
& \textbullet~Rhythm and pace parameters are reliable. \newline \textbullet~Joint velocity and event timing depend on \textbf{geometry}.
& \textbullet~Interpretation restricted to trustworthy components. \newline \textbullet~Acquisition geometry reported explicitly. \\
\addlinespace

Papanastasiou et al.\ \cite{papanastasiou2020}; Ni and Huang \cite{ni2020}
& \textbullet~Micro-Doppler identifies \textbf{individuals} at over 93\,\%.
& \textbullet~Gait is a \textbf{fingerprint}. \newline \textbullet~\textbf{Leave-one-subject-out} validation non-negotiable. \\
\addlinespace

Seyfio\u{g}lu et al.\ \cite{seyfioglu2018}; Park et al.\ \cite{park2016}
& \textbullet~Pretrained networks transfer well to new radar tasks.
& \textbullet~Motivates the \textbf{ResNet-18} transfer branch. \newline \textbullet~A frozen backbone suits a 58-subject cohort. \\
\addlinespace

Erol et al.\ \cite{erol2020}
& \textbullet~\textbf{ACGAN} synthesises realistic micro-Doppler. \newline \textbullet~Expands small training sets.
& \textbullet~Supports data augmentation. \newline \textbullet~Informs the augmentation recipe of Chapter~\ref{ch:methods}. \\
\addlinespace

Cai et al.\ \cite{cai2023}
& \textbullet~\textbf{STRIDE}: radar gait analysis for dementia-risk screening.
& \textbullet~Radar gait analysis generalises beyond Parkinson's. \\
\addlinespace

Meral and Ozbilgin \cite{neurosurgery2025}
& \textbullet~PD diagnosis from gait signals with deep learning. \newline \textbullet~Force sensors, not radar.
& \textbullet~PD gait carries a learnable signature. \newline \textbullet~\textbf{Radar} remained unexplored for subject-independent PD classification. \\

\bottomrule
\end{tabular}
\end{table}

\section{Research gap and positioning of this thesis}
\label{sec:gap}

The reviewed literature supports three conclusions.

\begin{enumerate}
  \item \textbf{What radar measures best is not what identifies the
        disease.} Radar micro-Doppler is a validated, information-rich way to
        record gait, but the quantities it measures most reliably (rhythm and
        pace) are not the ones most specific to Parkinson's (variability and
        absolute velocity) \cite{seifert2020,gurbuz2024}.

  \item \textbf{Deep learning on micro-Doppler is mature, but not for this
        task.} It performs well on broad problems such as activity recognition,
        person identification and young-versus-elderly discrimination, yet it has
        not been applied to subject-independent Parkinson's classification
        \cite{hayashi2021,papanastasiou2020,ni2020,lopezdelgado2026}.

  \item \textbf{Headline accuracies in this field are not trustworthy.} They are
        repeatedly undermined by limitations that recur from work to work: random
        hold-out splits that let the same person appear in training and test,
        shortcut learning on recording artifacts, cohorts of three to
        twenty-two subjects, an uncontrolled confound in \textbf{age}
        \cite{hayashi2021}, and, for Parkinson's disease, no radar classifier
        at all.
\end{enumerate}

\noindent Closing this gap is the purpose of this thesis, stated in
Chapter~\ref{ch:intro}, on the multi-node FMCW dataset of López-Delgado et al.\
\cite{lopezdelgado2026}. Methodologically, the work adopts together the
safeguards that the literature identifies as necessary but rarely combines:

\begin{itemize}
  \item \textbf{Leave-one-subject-out cross-validation} with subject-level score
        aggregation, to eliminate identity leakage.
  \item \textbf{Fixed-length windowing}, to suppress a trial-duration confound.
  \item \textbf{Explicit null floors}: deliberately simple reference
        predictors, recording duration alone and window count alone,
        reported beside every model to guard against the confounds
        highlighted by \cite{hayashi2021}.
  \item Subject-level \textbf{bootstrap confidence intervals}, for honest
        uncertainty reporting.
  \item \textbf{Gradient-based attention checks} (Grad-CAM), to verify that the
        network attends to the Doppler bands of the feet and the torso rather
        than to artifacts.
\end{itemize}

\noindent One safeguard the literature calls for could not be adopted: the
dataset records no ages, so the age confound is not controlled and is carried
as a limitation (Section~\ref{sec:limitations}). An uncontrolled confound can
only make a classifier look better than it is, so it cannot rescue a result
that already sits at the duration floor.

\noindent In doing so, the work aims to establish not merely whether Parkinson's
gait is detectable from radar, but whether it is detectable \emph{under
validation conditions strong enough to be trusted}.
